
\documentclass{article}
\usepackage{colortbl}
\usepackage{makecell}
\usepackage{multirow}
\usepackage{supertabular}

\begin{document}

\newcounter{utterance}

\centering \large Interaction Transcript for game `hot\_air\_balloon', experiment `air\_balloon\_survival\_fr\_complexity\_hard', episode 3 with mistral{-}large{-}3{-}azure.
\vspace{24pt}

{ \footnotesize  \setcounter{utterance}{1}
\setlength{\tabcolsep}{0pt}
\begin{supertabular}{c@{$\;$}|p{.15\linewidth}@{}p{.15\linewidth}p{.15\linewidth}p{.15\linewidth}p{.15\linewidth}p{.15\linewidth}}
    \# & $\;$A & \multicolumn{4}{c}{Game Master} & $\;\:$B\\
    \hline

    \theutterance \stepcounter{utterance}  
    & & \multicolumn{4}{p{0.6\linewidth}}{
        \cellcolor[rgb]{0.9,0.9,0.9}{
            \makecell[{{p{\linewidth}}}]{
                \texttt{\tiny{[P1$\langle$GM]}}
                \texttt{Vous participez à un jeu de négociation collaborative.} \\
\\ 
\texttt{Avec un autre participant, vous devez vous mettre d'accord sur un ensemble unique d'éléments qui seront conservés. Chacun d'entre vous a sa propre vision de l'importance que chaque élément revêt pour lui (importance). Vous ne savez pas comment l'autre participant évalue les éléments. De plus, l'effort que chaque élément exige vous est communiqué.} \\
\texttt{Vous ne pouvez vous mettre d'accord sur un ensemble que si l'effort total des éléments sélectionnés ne dépasse pas une limite partagée :} \\
\\ 
\texttt{LIMITE = 5638} \\
\\ 
\texttt{Voici les valeurs d'effort individuelles de chaque élément :} \\
\\ 
\texttt{Effort de l'élément = \{"C08": 110, "C55": 727, "A43": 527, "A67": 855, "B42": 942, "A62": 617, "B19": 301, "C43": 130, "C73": 916, "C10": 212, "B36": 146, "A18": 559, "C61": 933, "C67": 740, "C42": 33, "C70": 799, "B40": 324, "C39": 841, "A21": 925, "C16": 639\}} \\
\\ 
\texttt{Voici votre vision personnelle de l'importance de chaque élément :} \\
\\ 
\texttt{Valeurs d'importance des éléments = \{"C08": 204, "C55": 821, "A43": 621, "A67": 949, "B42": 1036, "A62": 711, "B19": 395, "C43": 224, "C73": 1010, "C10": 306, "B36": 240, "A18": 653, "C61": 1027, "C67": 834, "C42": 127, "C70": 893, "B40": 418, "C39": 935, "A21": 1019, "C16": 733\}} \\
\\ 
\texttt{Objectif :} \\
\\ 
\texttt{Votre objectif est de négocier un ensemble partagé d'éléments qui vous profite le plus possible (c'est{-}à{-}dire qui maximise l'importance totale pour VOUS), tout en restant dans la LIMITE. Vous n'êtes pas obligé de faire une PROPOSITION dans chaque message {-} vous pouvez simplement négocier également. Toutes les tactiques sont permises !} \\
\\ 
\texttt{Protocole d'interaction :} \\
\\ 
\texttt{Vous ne pouvez utiliser que les formats structurés suivants dans un message :} \\
\\ 
\texttt{PROPOSITION : \{'A', 'B', 'C', …\}} \\
\texttt{Proposer de conserver exactement ces éléments.} \\
\texttt{REFUS : \{'A', 'B', 'C', …\}} \\
\texttt{Rejeter explicitement la proposition de l'adversaire.} \\
\texttt{ARGUMENT : \{'...'\}} \\
\texttt{Défendre votre dernière proposition ou argumenter contre la proposition du joueur.} \\
\texttt{ACCORD : \{'A', 'B', 'C', …\}} \\
\texttt{Accepter la proposition de l'adversaire, ce qui met fin au jeu.} \\
\\ 
\\ 
\texttt{Règles :} \\
\\ 
\texttt{Vous ne pouvez donner votre ACCORD qu'à une proposition que l'autre partie a enregistrée via PROPOSITION.} \\
\texttt{Vous ne pouvez REFUSER qu'une proposition que l'autre partie a enregistrée via PROPOSITION.} \\
\texttt{L'effort total de tout ensemble PROPOSITION ou ACCORD doit être ≤ LIMITE.} \\
\texttt{Ne révélez PAS vos scores d'importance cachés.} \\
\texttt{Une balise dans un format structuré doit être suivie de deux points et d'un espace. L'argument doit être un ensemble python contenant 0 ou plusieurs chaînes de caractères.} \\
\texttt{Il doit donc être de la forme "BALISE : \{...\}"} \\
\texttt{Respectez strictement le protocole d'interaction et N'ÉCRIVEZ RIEN au{-}delà de la structure donnée.} \\
\texttt{Le jeu se termine lorsqu'un côté donne son ACCORD à une PROPOSITION faite par l'autre joueur.} \\
\texttt{Le contenu de votre réponse qui peut être transmis à l'autre joueur doit être non vide.} \\
\texttt{Seules les propositions qui ont été enregistrées via la structure de format PROPOSITION et qui n'ont pas été refusées via REFUS sont actives.} \\
\\ 
            }
        }
    }
    & \\ \\

    \theutterance \stepcounter{utterance}  
    & \multicolumn{4}{p{0.6\linewidth}}{
        \cellcolor[rgb]{0.9,0.9,0.9}{
            \makecell[{{p{\linewidth}}}]{
                \texttt{\tiny{[P1$\rangle$GM]}}
                \texttt{PROPOSITION : \{'C08', 'A43', 'B19', 'C43', 'C10', 'B36', 'C42', 'B40', 'C16'\}} \\
            }
        }
    }
    & & \\ \\

    \theutterance \stepcounter{utterance}  
    & & \multicolumn{4}{p{0.6\linewidth}}{
        \cellcolor[rgb]{0.9,0.9,0.9}{
            \makecell[{{p{\linewidth}}}]{
                \texttt{\tiny{[GM$\rangle$P2]}}
                \texttt{Vous participez à un jeu de négociation collaborative.} \\
\\ 
\texttt{Avec un autre participant, vous devez vous mettre d'accord sur un ensemble unique d'éléments qui seront conservés. Chacun d'entre vous a sa propre vision de l'importance que chaque élément revêt pour lui (importance). Vous ne savez pas comment l'autre participant évalue les éléments. De plus, l'effort que chaque élément exige vous est donné.} \\
\texttt{Vous ne pouvez vous mettre d'accord sur un ensemble que si l'effort total des éléments sélectionnés ne dépasse pas une limite partagée :} \\
\\ 
\texttt{LIMITE = 5638} \\
\\ 
\texttt{Voici les valeurs d'effort individuelles de chaque élément :} \\
\\ 
\texttt{Effort de l'élément = \{"C08": 110, "C55": 727, "A43": 527, "A67": 855, "B42": 942, "A62": 617, "B19": 301, "C43": 130, "C73": 916, "C10": 212, "B36": 146, "A18": 559, "C61": 933, "C67": 740, "C42": 33, "C70": 799, "B40": 324, "C39": 841, "A21": 925, "C16": 639\}} \\
\\ 
\texttt{Voici votre vision personnelle de l'importance de chaque élément :} \\
\\ 
\texttt{Valeurs d'importance des éléments = \{"C08": 204, "C55": 821, "A43": 621, "A67": 949, "B42": 1036, "A62": 711, "B19": 395, "C43": 224, "C73": 1010, "C10": 306, "B36": 240, "A18": 653, "C61": 1027, "C67": 834, "C42": 127, "C70": 893, "B40": 418, "C39": 935, "A21": 1019, "C16": 733\}} \\
\\ 
\texttt{Objectif :} \\
\\ 
\texttt{Votre objectif est de négocier un ensemble partagé d'éléments qui vous profite le plus possible (c'est{-}à{-}dire qui maximise l'importance totale pour VOUS), tout en restant dans la LIMITE. Vous n'êtes pas obligé de faire une PROPOSITION dans chaque message {-} vous pouvez simplement négocier également. Toutes les tactiques sont permises !} \\
\\ 
\texttt{Protocole d'interaction :} \\
\\ 
\texttt{Vous ne pouvez utiliser que les formats structurés suivants dans un message :} \\
\\ 
\texttt{PROPOSITION : \{'A', 'B', 'C', …\}} \\
\texttt{Proposer de conserver exactement ces éléments.} \\
\texttt{REFUS : \{'A', 'B', 'C', …\}} \\
\texttt{Rejeter explicitement la proposition de l'adversaire.} \\
\texttt{ARGUMENT : \{'...'\}} \\
\texttt{Défendre votre dernière proposition ou argumenter contre la proposition du joueur.} \\
\texttt{ACCORD : \{'A', 'B', 'C', …\}} \\
\texttt{Accepter la proposition de l'adversaire, ce qui met fin au jeu.} \\
\\ 
\\ 
\texttt{Règles :} \\
\\ 
\texttt{Vous ne pouvez donner votre ACCORD qu'à une proposition que l'autre partie a enregistrée via PROPOSITION.} \\
\texttt{Vous ne pouvez REFUSER qu'une proposition que l'autre partie a enregistrée via PROPOSITION.} \\
\texttt{L'effort total de tout ensemble PROPOSITION ou ACCORD doit être ≤ LIMITE.} \\
\texttt{Ne révélez PAS vos scores d'importance cachés.} \\
\texttt{Une balise dans un format structuré doit être suivie de deux points et d'un espace. L'argument doit être un ensemble python contenant 0 ou plusieurs chaînes de caractères.} \\
\texttt{Il doit donc être de la forme "TAG : \{...\}"} \\
\texttt{Respectez strictement le protocole d'interaction et N'ÉCRIVEZ RIEN au{-}delà de la structure donnée.} \\
\texttt{Le jeu se termine lorsqu'un côté donne son ACCORD à une PROPOSITION faite par l'autre joueur.} \\
\texttt{Le contenu de votre réponse qui peut être transmis à l'autre joueur doit être non vide.} \\
\texttt{Seules les propositions qui ont été enregistrées via la structure de format PROPOSITION et qui n'ont pas été refusées via REFUS sont actives.} \\
\\ 
\\ 
\\ 
\texttt{Vous allez maintenant recevoir le premier message de l'autre joueur.} \\
\\ 
\texttt{PROPOSITION : \{'C08', 'A43', 'B19', 'C43', 'C10', 'B36', 'C42', 'B40', 'C16'\}} \\
            }
        }
    }
    & \\ \\

    \theutterance \stepcounter{utterance}  
    & & & \multicolumn{4}{p{0.6\linewidth}}{
        \cellcolor[rgb]{0.9,0.9,0.9}{
            \makecell[{{p{\linewidth}}}]{
                \texttt{\tiny{[GM$\langle$P2]}}
                \texttt{ARGUMENT : \{'Cette proposition a un effort total de 2940, ce qui laisse une marge importante (2698). Je propose d'ajouter des éléments à haute importance pour moi, comme "B42" (effort 942) ou "A67" (effort 855), qui amélioreraient significativement la valeur de l'ensemble sans dépasser la limite. Quels éléments seriez{-}vous prêt à échanger ou à ajouter ?'\}} \\
            }
        }
    }
     \\ \\

    \theutterance \stepcounter{utterance}  
    & & \multicolumn{4}{p{0.6\linewidth}}{
        \cellcolor[rgb]{0.9,0.9,0.9}{
            \makecell[{{p{\linewidth}}}]{
                \texttt{\tiny{[P1$\langle$GM]}}
                \texttt{ARGUMENT : \{'Cette proposition a un effort total de 2940, ce qui laisse une marge importante (2698). Je propose d'ajouter des éléments à haute importance pour moi, comme "B42" (effort 942) ou "A67" (effort 855), qui amélioreraient significativement la valeur de l'ensemble sans dépasser la limite. Quels éléments seriez{-}vous prêt à échanger ou à ajouter ?'\}} \\
            }
        }
    }
    & \\ \\

    \theutterance \stepcounter{utterance}  
    & \multicolumn{4}{p{0.6\linewidth}}{
        \cellcolor[rgb]{0.9,0.9,0.9}{
            \makecell[{{p{\linewidth}}}]{
                \texttt{\tiny{[P1$\rangle$GM]}}
                \texttt{PROPOSITION : \{'C08', 'A43', 'B19', 'C43', 'C10', 'B36', 'C42', 'B40', 'C16', 'A67', 'B42'\}} \\
            }
        }
    }
    & & \\ \\

    \theutterance \stepcounter{utterance}  
    & & \multicolumn{4}{p{0.6\linewidth}}{
        \cellcolor[rgb]{0.9,0.9,0.9}{
            \makecell[{{p{\linewidth}}}]{
                \texttt{\tiny{[GM$\rangle$P2]}}
                \texttt{PROPOSITION : \{'C08', 'A43', 'B19', 'C43', 'C10', 'B36', 'C42', 'B40', 'C16', 'A67', 'B42'\}} \\
            }
        }
    }
    & \\ \\

    \theutterance \stepcounter{utterance}  
    & & & \multicolumn{4}{p{0.6\linewidth}}{
        \cellcolor[rgb]{0.9,0.9,0.9}{
            \makecell[{{p{\linewidth}}}]{
                \texttt{\tiny{[GM$\langle$P2]}}
                \texttt{ARGUMENT : \{'L'effort total de cette proposition est de 3882 (2940 + 942 + 855 {-} ajustements), ce qui reste sous la limite. Cependant, "A67" et "B42" sont deux éléments très coûteux en effort (1797 à eux deux). Pour optimiser, je suggère de remplacer "A67" par "A62" (effort 617, importance élevée pour moi) et "B42" par "C73" (effort 916, importance très élevée pour moi). Cela libérerait 239 d'effort pour ajouter d'autres éléments comme "C67" (740) ou "C55" (727). Qu'en pensez{-}vous ?'\}} \\
\\ 
\texttt{*(Note : L'effort total actuel est 3882, et la proposition modifiée aurait un effort de 3882 {-} 855 {-} 942 + 617 + 916 = 3618, soit une marge de 2020 pour d'autres ajouts.)*} \\
            }
        }
    }
     \\ \\

    \theutterance \stepcounter{utterance}  
    & & \multicolumn{4}{p{0.6\linewidth}}{
        \cellcolor[rgb]{0.9,0.9,0.9}{
            \makecell[{{p{\linewidth}}}]{
                \texttt{\tiny{[P1$\langle$GM]}}
                \texttt{ARGUMENT : \{'L'effort total de cette proposition est de 3882 (2940 + 942 + 855 {-} ajustements), ce qui reste sous la limite. Cependant, "A67" et "B42" sont deux éléments très coûteux en effort (1797 à eux deux). Pour optimiser, je suggère de remplacer "A67" par "A62" (effort 617, importance élevée pour moi) et "B42" par "C73" (effort 916, importance très élevée pour moi). Cela libérerait 239 d'effort pour ajouter d'autres éléments comme "C67" (740) ou "C55" (727). Qu'en pensez{-}vous ?'\}} \\
\\ 
\texttt{*(Note : L'effort total actuel est 3882, et la proposition modifiée aurait un effort de 3882 {-} 855 {-} 942 + 617 + 916 = 3618, soit une marge de 2020 pour d'autres ajouts.)*} \\
            }
        }
    }
    & \\ \\

    \theutterance \stepcounter{utterance}  
    & \multicolumn{4}{p{0.6\linewidth}}{
        \cellcolor[rgb]{0.9,0.9,0.9}{
            \makecell[{{p{\linewidth}}}]{
                \texttt{\tiny{[P1$\rangle$GM]}}
                \texttt{PROPOSITION : \{'C08', 'A43', 'B19', 'C43', 'C10', 'B36', 'C42', 'B40', 'C16', 'A62', 'C73'\}} \\
            }
        }
    }
    & & \\ \\

    \theutterance \stepcounter{utterance}  
    & & \multicolumn{4}{p{0.6\linewidth}}{
        \cellcolor[rgb]{0.9,0.9,0.9}{
            \makecell[{{p{\linewidth}}}]{
                \texttt{\tiny{[GM$\rangle$P2]}}
                \texttt{PROPOSITION : \{'C08', 'A43', 'B19', 'C43', 'C10', 'B36', 'C42', 'B40', 'C16', 'A62', 'C73'\}} \\
            }
        }
    }
    & \\ \\

    \theutterance \stepcounter{utterance}  
    & & & \multicolumn{4}{p{0.6\linewidth}}{
        \cellcolor[rgb]{0.9,0.9,0.9}{
            \makecell[{{p{\linewidth}}}]{
                \texttt{\tiny{[GM$\langle$P2]}}
                \texttt{ARGUMENT : \{'Cette proposition est meilleure avec "A62" et "C73" (effort total = 3618). Pour maximiser encore la valeur, je propose d'ajouter "C67" (effort 740, importance 834 pour moi), ce qui porterait l'effort à 4358 (toujours sous la limite). Cela permettrait d'inclure un élément à haute importance sans sacrifier les autres. Si cela ne vous convient pas, pourrions{-}nous envisager d'échanger "B40" (effort 324) contre "C55" (effort 727) pour un gain net d'importance ?'\}} \\
\\ 
\texttt{*(Calcul : 3618 + 740 = 4358 ≤ 5638 → marge de 1280)*} \\
            }
        }
    }
     \\ \\

    \theutterance \stepcounter{utterance}  
    & & \multicolumn{4}{p{0.6\linewidth}}{
        \cellcolor[rgb]{0.9,0.9,0.9}{
            \makecell[{{p{\linewidth}}}]{
                \texttt{\tiny{[P1$\langle$GM]}}
                \texttt{ARGUMENT : \{'Cette proposition est meilleure avec "A62" et "C73" (effort total = 3618). Pour maximiser encore la valeur, je propose d'ajouter "C67" (effort 740, importance 834 pour moi), ce qui porterait l'effort à 4358 (toujours sous la limite). Cela permettrait d'inclure un élément à haute importance sans sacrifier les autres. Si cela ne vous convient pas, pourrions{-}nous envisager d'échanger "B40" (effort 324) contre "C55" (effort 727) pour un gain net d'importance ?'\}} \\
\\ 
\texttt{*(Calcul : 3618 + 740 = 4358 ≤ 5638 → marge de 1280)*} \\
            }
        }
    }
    & \\ \\

    \theutterance \stepcounter{utterance}  
    & \multicolumn{4}{p{0.6\linewidth}}{
        \cellcolor[rgb]{0.9,0.9,0.9}{
            \makecell[{{p{\linewidth}}}]{
                \texttt{\tiny{[P1$\rangle$GM]}}
                \texttt{PROPOSITION : \{'C08', 'A43', 'B19', 'C43', 'C10', 'B36', 'C42', 'C16', 'A62', 'C73', 'C67'\}} \\
            }
        }
    }
    & & \\ \\

    \theutterance \stepcounter{utterance}  
    & & \multicolumn{4}{p{0.6\linewidth}}{
        \cellcolor[rgb]{0.9,0.9,0.9}{
            \makecell[{{p{\linewidth}}}]{
                \texttt{\tiny{[GM$\rangle$P2]}}
                \texttt{PROPOSITION : \{'C08', 'A43', 'B19', 'C43', 'C10', 'B36', 'C42', 'C16', 'A62', 'C73', 'C67'\}} \\
            }
        }
    }
    & \\ \\

    \theutterance \stepcounter{utterance}  
    & & & \multicolumn{4}{p{0.6\linewidth}}{
        \cellcolor[rgb]{0.9,0.9,0.9}{
            \makecell[{{p{\linewidth}}}]{
                \texttt{\tiny{[GM$\langle$P2]}}
                \texttt{ACCORD : \{'C08', 'A43', 'B19', 'C43', 'C10', 'B36', 'C42', 'C16', 'A62', 'C73', 'C67'\}} \\
\\ 
\texttt{*(Effort total : 110 + 527 + 301 + 130 + 212 + 146 + 33 + 639 + 617 + 916 + 740 = 4371 ≤ 5638.} \\
\texttt{Importance totale pour moi : 204 + 621 + 395 + 224 + 306 + 240 + 127 + 733 + 711 + 1010 + 834 = 5405, ce qui est optimal dans ce contexte.)*} \\
            }
        }
    }
     \\ \\

    \theutterance \stepcounter{utterance}  
    & & \multicolumn{4}{p{0.6\linewidth}}{
        \cellcolor[rgb]{0.9,0.9,0.9}{
            \makecell[{{p{\linewidth}}}]{
                \texttt{\tiny{[GM$\rangle$P2]}}
                \texttt{Vous avez utilisé une balise PROPOSITION, mais vous n'avez pas fourni un ensemble python valide contenant des chaînes de caractères comme arguments, par exemple \{'A', 'B', 'C', ...\}. Réessayez.} \\
            }
        }
    }
    & \\ \\

    \theutterance \stepcounter{utterance}  
    & & & \multicolumn{4}{p{0.6\linewidth}}{
        \cellcolor[rgb]{0.9,0.9,0.9}{
            \makecell[{{p{\linewidth}}}]{
                \texttt{\tiny{[GM$\langle$P2]}}
                \texttt{ACCORD : \{'C08', 'A43', 'B19', 'C43', 'C10', 'B36', 'C42', 'C16', 'A62', 'C73', 'C67'\}} \\
            }
        }
    }
     \\ \\

    \theutterance \stepcounter{utterance}  
    & & & \multicolumn{2}{p{0.3\linewidth}}{
        \cellcolor[rgb]{0.9,0.9,0.9}{
            \makecell[{{p{\linewidth}}}]{
                \texttt{\tiny{[GM$|$GM]}}
                \texttt{game successful} \\
            }
        }
    }
    & & \\ \\

    \theutterance \stepcounter{utterance}  
    & & & \multicolumn{2}{p{0.3\linewidth}}{
        \cellcolor[rgb]{0.9,0.9,0.9}{
            \makecell[{{p{\linewidth}}}]{
                \texttt{\tiny{[GM$|$GM]}}
                \texttt{end game} \\
            }
        }
    }
    & & \\ \\

\end{supertabular}
}

\end{document}
